Machine learning algorithms and their training datasets have grown substantially in both size and complexity over the past decade. This increased model complexity has made it challenging to interpret and predict their behavior in unobserved scenarios. Hence, many applications that involve societal decisions still rely on simple, interpretable models like linear regression, often after feature engineering. Examples of such applications include predicting national housing prices, estimating wages across industries, forecasting loan amounts across banks, predicting life insurance premiums across groups, and projecting energy consumption across communities \citep{cohen2024fairness}. 

A shared safety and sometimes legal concern across the above applications is the potential for wildly different model qualities for different distributions, i.e., outputting a notably worse model for some source data distributions~\citep{data2014seizing, barocas2016big,hardt2016equality,veale2018fairness,selbst2019fairness,berk2021fairness, corbett2023measure,chouldechova2016fair,kleinberg2018algorithmic, agarwal2019fair,cohen2024fairness,svwz24}. Specifically, consider fitting a linear model $\vx\in\R^d$ to make real predictions on some task over $m$ groups where group $i$'s dataset consists of $n_i$ entries and is denoted by $S_i=\{(\va_i^j, b_i^j)\}_{j\in [n_i]}$. The \textit{utilitarian} or the total-cost-minimizing objective minimizes the average squared prediction error across groups, i.e.,        
\begin{align}
    \min_{\vx\in\R^d}\frac{1}{m}\sum_{i\in[m]}\frac{1}{n_i}\norm{\mA_{S_i}\vx-\vb_{S_i}}_2^2\enspace,\label{eq:utilitarian}
\end{align}
where $\mA_{S_i} \coloneqq [\va_i^1 \dots \va_i^{n_i}]^{\top}\in \R^{n_i \times d}$ is the feature matrix and $\vb_{S_i} \coloneqq [b_i^1 \dots b_i^{n_i}]^{\top} \in \R^{n_i}$ is the label vector for group $i\in[m]$. 

Due to the inherent heterogeneity of the datasets, the model derived by optimizing the objective \eqref{eq:utilitarian} may be particularly detrimental to some groups, as the prediction error may be disproportionately higher for these groups. To overcome these limitations, the following \textit{egalitarian} or group Distributionally Robust Optimization (DRO) objective  has been considered in several recent works \citep{ben2013robust, duchi2016statistics, sagawa2019distributionally, levy2020large, soma2022optimal,aakmrz22,svwz24},
\begin{align}
    \min_{\vx\in\R^d} \max_{i\in[m]}\frac{1}{n_i}\norm{\mA_{S_i}\vx-\vb_{S_i}}_2^2\enspace.\label{eq:main_objective}
\end{align}
Objective \ref{eq:main_objective} is the ``fairest'' objective among all objectives that balance utility and distributional robustness~\citep{kleinberg2018algorithmic,chouldechova2018frontiers,asadpour2022sequential,chen2022fair,rahmattalabi2019exploring,golrezaei2024online}. Since objective \ref{eq:main_objective} is a convex problem, it is natural to apply standard black-box optimization techniques to solve it. However, we identify several challenges in applying existing methods:
    \paragraph{Efficient first-order algorithms have geometry-dependent rates.} To our knowledge, using an efficient first-order method (such as sub-gradient descent) will incur a geometry-dependent runtime. In particular, if the matrices $\mA_{S_i}$ or if the stacked matrix $\mA\coloneqq [\mA_{S_1}^{\top} \dots \mA_{S_M}^{\top}]^{\top}$ are poorly conditioned, then this will be reflected accordingly in the convergence rates. This is a drawback of the existing results by \citet{aakmrz22} and \citet{svwz24}.
    \paragraph{Objective \eqref{eq:main_objective} is not smooth.} Since the objective is the pointwise maximum of several continuous functions, the derivative is not well-defined at the points at which the maximizing function changes. Thus, applying subgradient descent to this objective without a tailored analysis will yield a rather unimpressive $1/\eps^{2}$ dependence in the iteration complexity. 
    \paragraph{Min-max optimization/regret minimization approaches have a $1/\eps^{2}$ dependence on iteration complexity.} Since problem \ref{eq:main_objective} is a min-max optimization objective, it is also natural to try to use game theory-inspired approaches that use some oracle (such as gradients) for each group as a black box. For instance, we can cast objective \ref{eq:main_objective} as a repeated game between a min player (equipped with a no-regret algorithm) and a max player (equipped with the best response oracle). The main shortcoming of this approach is that even though the function for each group is smooth, the iteration complexity (to get $\eps$ average regret) for smooth online convex optimization still has an unimpressive $1/\eps^2$ dependence (as opposed to $1/\eps$ for smooth convex optimization)~\citep{soma2022optimal, zhang2024stochastic}. Thus, this approach is no better than optimizing \eqref{eq:main_objective} using sub-gradient descent.
    \paragraph{Interior point methods have a poor iteration complexity for large $m$.} Another natural approach (that can partially address the previous two issues), following the discussion by \citet[Section 6.4]{boyd2004convex}, is to rewrite problem \ref{eq:main_objective} in its epigraph form and use an interior point method (IPM) to solve the resulting problem (which, in this case, is a quadratically constrained linear program). Unfortunately, this will give an algorithm whose analysis is only known to yield an iteration complexity of $O(\sqrt{m})$, where each iteration solves a linear system in matrices of the form $\mA^{\top}\mB\mA$ for a block-diagonal $\mB$ (see \Cref{rem:linear_systems}).
    A na\"ive implementation of this algorithm will thus have a superlinear runtime in the number of groups, which is undesirable when the number of groups is large. Alternately, consider an example in which we copy each group $k$ times in the objective. The new objective value does not change from the original objective value, but the iteration complexity from the IPM now blows up to $\sqrt{mk}$. This also signals to us that we should search for an algorithm whose iteration complexity is mostly independent from $m$. %Furthermore, note that \eqref{eq:main_objective} is not a linear program when at least one group $i$ is such that $n_i > 1$. So, we cannot immediately apply recent advances in linear programming that get iteration complexities independent of the number of constraints \citep{ls19}.

Hence, designing an algorithm without these shortcomings requires novel ideas. 

\subsection{Our Results}
\label{sec:gp_our_results}

In this paper, we present a new algorithm (\Cref{alg:fair_regression}) to approximately optimize objective \ref{eq:main_objective}, which addresses the aforementioned difficulties. We state the iteration complexity of our algorithm in the following theorem. 

\begin{mainthm}[Robust regression]\label{mainthm:fair_regression_iteration_complexity}
Let $\mA_{S_i} \in \R^{n_i \times d}$ and $\vb_{S_i} \in \R^{n_i}$ for all $i\in[m]$. Denote their concatenations by $\mA\coloneqq [\mA_{S_1}^{\top} \dots \mA_{S_M}^{\top}]^{\top}\in\R^{n\times d}$ and $\vb \coloneqq [\vb_{S_1}^{\top}\dots\vb_{S_M}^{\top}]^{\top}\in \R^n$ where $n:= \sum_{i\in[m]}n_i$. Let $\eps > 0$. Then \Cref{alg:fair_regression} returns $\xhat$ such that,
\begin{align}
    \max_{i\in[m]} \frac{1}{\sqrt{n_i}}\norm{\mA_{S_i}\xhat-\vb_{S_i}}_2 \le (1+\eps)\cdot\min_{\vx\in\R^d}\max_{i\in[m]} \frac{1}{\sqrt{n_i}}\norm{\mA_{S_i}\vx-\vb_{S_i}}_2\enspace,\label{eq:objective_approx}
\end{align}
and it runs in
\begin{align*}
    O\inparen{\frac{\min\inbraces{\mathsf{rank}(\mA),m}^{1/3}\inparen{\logv{\frac{n\log m}{\eps}}^{14/3}+\logv{m}}}{\eps^{2/3}}}
\end{align*}
linear-system-solves in matrices of the form $\mA^{\top}\mB\mA$, where $\mB$ is a block-diagonal matrix for which block $i$ has size $n_i \times n_i$.
\end{mainthm}
We prove \Cref{mainthm:fair_regression_iteration_complexity} in \Cref{sec:gp_robust_proof}. We compare the guarantee of \Cref{mainthm:fair_regression_iteration_complexity} against the other baselines in \Cref{table:compare}. 

\def\arraystretch{1.5}
\begin{table}[h]
\centering
\begin{tabular}{l c c} 
\toprule\toprule
\textbf{Algorithm}                                                                                                     & \textbf{Iteration Complexity}                                                                              & \textbf{Each Iteration}                     \\ \midrule\midrule
\makecell[l]{Subgradient descent}                                                                                           & \makecell[c]{$\frac{\norm{\xstar}_2 \max_{1 \le i \le m} \tfrac{1}{\sqrt{n_i}}\opnorm{\mA_{S_i}}}{\eps^{2}}$}                                   & \makecell[c]{Evaluate $\nabla f(\vx)$}                    \\ \midrule
\makecell[l]{Nesterov acceleration\\ on smoothened objective}                                                                 & \makecell[c]{$\frac{\norm{\xstar}_2 \left(\max_{1 \le i \le m} \tfrac{1}{\sqrt{n_i}}\opnorm{\mA_{S_i}}\right)^{1/2}}{\eps}$}                      & \makecell[c]{Evaluate $\nabla \ftilde_{\beta,\delta}(\vx)$} \\ \midrule
\makecell[l]{\cite{aakmrz22}}                                                                                               & \makecell[c]{$\frac{\norm{\xstar}_2 \max_{1 \le i \le m} \tfrac{1}{\sqrt{n_i}}\opnorm{\mA_{S_i}}}{\eps}$}                                       & \makecell[c]{Evaluate $\nabla \ftilde_{\beta,\delta}(\vx)$} \\ \midrule
\makecell[l]{Interior point with log barrier\\ \citep{boyd2004convex}}                                                         & \makecell[c]{$m^{1/2} \log\left(\frac{1}{\eps}\right)$}                                                                    & \makecell[c]{Linear-system-solve\\ in $\mA^{\top}\mB\mA$}    \\ \midrule
\makecell[l]{\textbf{This paper} \\ (na\"ive geometry)}                                                                                          & \makecell[c]{$\frac{m^{1/3}}{\eps^{2/3}}$}                                                                                & \makecell[c]{Linear-system-solve\\ in $\mA^{\top}\mB\mA$}    \\ \midrule
\makecell[l]{$\ell_{\infty}$ regression with Lewis \\weights \citep{jls21}}                                                    & \makecell[c]{$\frac{\rank{\mA}^{1/3}}{\eps^{2/3}}$}                                                                       & \makecell[c]{Linear-system-solve\\ in $\mA^{\top}\mD\mA$}    \\ \midrule
\makecell[l]{$\ell_{\infty}$ regression with IPM\\ \citep{ls19}}                                                               & \makecell[c]{$\rank{\mA}^{1/2} \log\left(\frac{1}{\eps}\right)$}                                                           & \makecell[c]{Linear-system-solve\\ in $\mA^{\top}\mD\mA$}    \\ \midrule
\makecell[l]{\textbf{This paper (\Cref{mainthm:fair_regression_iteration_complexity})}}                                                                                                    & \makecell[c]{$\frac{\min\inbraces{\rank{\mA},m}^{1/3}}{\eps^{2/3}}$}                                                               & \makecell[c]{Linear-system-solve\\ in $\mA^{\top}\mB\mA$}    \\ \bottomrule\bottomrule
\end{tabular}
\caption{The complexities of algorithms for optimizing \eqref{eq:main_objective} or for the special case of $\ell_{\infty}$ regression, assuming $\opt=1$ (the first three guarantees are additive approximations) and ignoring $\mathrm{polylog}(n,m)$ terms. We write $\mD$ to be a diagonal matrix and $\mB$ to be a block-diagonal matrix where each block has size $(n_i + o(1)) \times (n_i + o(1))$. We remark that in the special case where $n_i=1$, our algorithm exactly recovers guarantees of \cite{jls21}. We stress that we include the references to $\ell_{\infty}$ regression only to show that our algorithm is no worse than that of \cite{jls21} in this special case of $n_i=1$ for all $i$, and none of their algorithms apply to our general setting.\label{table:compare}}
\end{table}

Unlike the aforementioned first-order methods, our algorithm has no geometry-dependent terms. Additionally, our algorithm improves over the standard log-barrier IPM when the desired accuracy $\eps \ge m^{-1/4}$ --- this improvement is more pronounced when $m \gg \rank{\mA}$, i.e. when the number of data sources is much larger than the dimension of the parameter vector $\vx$. Additionally, for $\eps \ge \rank{\mA}^{-1/4}$, our guarantee matches the best known guarantee for $\ell_{\infty}$ regression \citep{ls19,jls21}.

\begin{remark}[Why use linear-system-solve complexity?]\label{rem:linear_systems}
    We benchmark our algorithms using the number of linear-system-solves for a few reasons. First, this is typically how second-order algorithms are compared, such as interior point methods for linear programming \citep{ls19}. Second, the particular structure of the linear-system-solves presents the possibility of a faster amortized runtime for the systems over the algorithm's run. This observation, combined with an understanding of how the linear systems changed between iterations, was used recently used to achieve fast runtimes for linear programming \citep{ls19} and $\ell_{\infty}$ regression \citep{ajk24}.
\end{remark}

\paragraph{Interpolating between robust and nonrobust optimization.} We also study the following family of objectives that interpolate between objectives \ref{eq:utilitarian} and \ref{eq:main_objective} for different values of $p\geq 2$,
\begin{align}
    \min_{\vx\in\R^d}\frac{1}{m}\sum_{i\in[m]}\left(\frac{1}{n_i}\norm{\mA_{S_i}\vx-\vb_{S_i}}_2^2\right)^{p/2}\enspace.\label{eq:interpolation}
\end{align}
In particular, note that choosing $p=2$ in the above objective gives us the average least-squares problem in objective \ref{eq:utilitarian}, while $p\to\infty$ recovers objective \ref{eq:main_objective}. Varying $p$ from $2$ to $\infty$ and minimizing gives solutions that interpolate between utilitarian and egalitarian approaches, allowing for a smooth trade-off between utility and robustness. To this end, we give \Cref{alg:gp_regression_alg} to approximately optimize objective \ref{eq:interpolation} and prove the following guarantee about its iteration complexity.

\begin{mainthm}[Trading off utility and robustness]\label{mainthm:gp_regression_iteration_complexity}
Let $\mA_{S_i} \in \R^{n_i \times d}$ and $\vb_{S_i} \in \R^{n_i}$ for all $i\in[m]$. Denote their concatenations by $\mA\coloneqq [\mA_{S_1}^{\top} \dots \mA_{S_M}^{\top}]^{\top}\in\R^{n\times d}$ and $\vb \coloneqq [\vb_{S_1}^{\top}\dots\vb_{S_M}^{\top}]^{\top}\in \R^n$ where $n \coloneqq \sum_{i\in[m]}n_i$. Let $p \ge 2$ and $\eps > 0$. Then \Cref{alg:gp_regression_alg} returns $\xhat$ such that,
\begin{align}\label{eq:interpolation_approx}
    \inparen{\sum_{i=1}^m \inparen{\frac{1}{\sqrt{n_i}}\norm{\mA_{S_i}\xhat-\vb_{S_i}}_2}^p}^{1/p} \le (1+\eps)\cdot\min_{\vx\in\R^d}\inparen{\sum_{i=1}^m \inparen{\frac{1}{\sqrt{n_i}}\norm{\mA_{S_i}\vx-\vb_{S_i}}_2}^p}^{1/p}
\end{align}
and runs in
\begin{align*}          
    O\inparen{p^{O(1)}\min\inbraces{\rank{\mA},m}^{\frac{p-2}{3p-2}}\logv{\frac{pd}{\eps}}^3}
\end{align*}
linear-system-solves in matrices of the form $\mA^{\top}\mB\mA$, where $\mB$ is a block-diagonal matrix for which block $i$ has size $n_i \times n_i$.
\end{mainthm}
We prove \Cref{mainthm:gp_regression_iteration_complexity} in \Cref{sec:gp_interpolating_proof}.

In the special case where $n_i=1$ for all $i$ (and therefore the problem is $\ell_p$ regression for $p \ge 2$), the complexity promised by \Cref{mainthm:gp_regression_iteration_complexity} is comparable to that promised by \citet{jls21} for $\ell_p$ regression. The main difference is that our iteration complexity is unconditionally polynomial in $p$. In contrast, the comparable result from \cite{jls21} seems to require mild assumptions on the problem parameters (see the ``Discussion on numerical stability'' by \citet[Section 4]{jls21}).

\begin{remark}[Large values of $p$] 
    Note that for values of $p$ larger than $\log(m)$, solving \eqref{eq:main_objective} is almost equivalent to solving \eqref{eq:interpolation}. To intuitively see this, first recall that for any vector $\vx\in\R^d$ and $p=\log_2(m)$ we have, $\|\vx\|_\infty \leq \|\vx\|_{p} \leq 2\cdot\|\vx\|_\infty$. This implies that for all $i\in[m]$ we have the following for objective \eqref{eq:interpolation} (for $p=\log_2(m)$) for any $\vx\in\R^d$,
    \begin{align*}
        \max_{i\in[m]}\norm{\mA_{S_i}\vx-\vb_{S_i}}_2 \leq \left(\sum_{i\in[m]}\norm{\mA_{S_i}\vx-\vb_{S_i}}_2^{p}\right)^{1/p}\leq 2\cdot\max_{i\in[m]}\norm{\mA_{S_i}\vx-\vb_{S_i}}_2\enspace. 
    \end{align*}
    In particular, this means that minimizing the interpolating objective \eqref{eq:interpolation} also minimizes the robust objective \eqref{eq:main_objective} (up to numerical constants) and vice versa. Thus, for $p=\Omega\left(\log_2(m)\right)$, for our intended applications, it makes sense to minimize the robust objective instead. This is why, in \Cref{mainthm:gp_regression_iteration_complexity}, we do not care too much about the exponent on $p$ in the iteration complexity. Our main goal is to show that we can get a $O(\mathsf{poly}(p,\logv{\frac{1}{\eps}})\min\inbraces{\rank{\mA},m}^{1/3})$ iteration complexity.
\end{remark}

\subsection{Prior Results, Connections, and Open Problems}
Here, we discuss prior work that conceptually and technically relates to ours. We then suggest natural directions for future work.

\paragraph{Multi-distribution learning.}
Many learning problems involve multiple data sources, for instance, when multiple agents generate their data independently. One can formulate these multi-distribution problems as standard learning/optimization problems by considering a mixture of their distributions, as in objective \ref{eq:utilitarian}. However, this approach often biases solutions toward dominant data sources, leading to poor performance on outliers—an issue stemming from statistical heterogeneity. This limitation motivates the study of multi-objective optimization problems~\citep{miettinen1999nonlinear, ehrgott2005multicriteria}, where each agent \( m \) has a distribution \( \mathcal{D}_m \) that defines its objective as \( \mathbb{E}_{z\sim \mathcal{D}_m}[f(\vx_m;z)] \), and where models \( \vx_m \) can vary across agents---a framework known as personalization. 

One of the earliest algorithms for such problems was introduced by \citet{blum2017collaborative}, where each agent's objective must be minimized to a pre-specified threshold \( \epsilon \) with high probability, framed within a PAC learning framework~\citep{valiant1984theory,vapnik2013nature}. Subsequent research has refined these algorithms, achieving optimal sample complexity guarantees for learning from multiple distributions~\citep{chen2018tight,nguyen2018improved,hanneke2019value,haghtalab2022demand,zhang2024optimal}. Our objectives \( \ref{eq:main_objective} \) and \( \ref{eq:interpolation} \) offer different approaches to multi-distribution learning, where data distributions correspond to empirical agent distributions. In particular, \citet{mohri2019agnostic} analyzed objective \( \ref{eq:main_objective} \) to establish generalization bounds for unknown mixtures of agents' distributions. 

Beyond sample efficiency, researchers have also examined other challenges, such as communication costs in large-scale distributed optimization~\citep{mcmahan2016s}. A particularly relevant study is that of \citet{bullins2021stochastic}, which employs an efficient distributed quadratic sub-solver~\citep{woodworth2020local,patel2024limits} to implement an inexact Newton method for optimizing quasi-self-concordant functions (see \Cref{defn:qsc}).

\paragraph{Group fairness.} 
Recently, interest in algorithmic fairness has intensified~\citep{barocas2016big, abebe2020roles, kasy2021fairness} with researchers exploring fairness across various domains, including supervised learning~\citep{calders2009building, dwork2012fairness, hardt2016equality, kusner2017counterfactual, goel2018non, ustun2019fairness}, resource allocation~\citep{bertsimas2011price, bertsimas2012efficiency, hooker2012combining, donahue2020fairness, manshadi2021fair}, scheduling~\citep{mulvany2021fair}, online matching~\citep{chierichetti2019matroids, ma2023fairness}, assortment planning~\citep{singh2018fairness, biega2018equity, singh2019policy, chen2022fair}, and facility location~\citep{gupta2022socially}. The extensive literature on algorithmic fairness falls into three main categories: (1) individual fairness, which ensures that similar individuals receive comparable predictions~\citep{dwork2012fairness, loi2019philosophical, chen2022fair}, (2) group fairness, which aims for equal treatment of different demographic groups, often in terms of resource allocation or performance parity~\citep{singh2018fairness, balseiro2021regularized}, and (3) subgroup fairness, which blends aspects of both individual and group fairness~\citep{kearns2018preventing, kearns2019empirical}. 

This paper focuses on a well-studied group fairness notion in machine learning literature: the group DRO problem~\citep{ben2013robust, duchi2016statistics, sagawa2019distributionally}. The idea of interpolating between robustness and utility is also common~\citep{golrezaei2024online} and closely related to multi-objective optimization, where scalarization~\citep{miettinen1999nonlinear, ehrgott2005multicriteria} helps recover desired solutions along the Pareto frontier. 

\paragraph{Linear programming and $\ell_p$ regression.} 
In the last several years, there has been a surge of work in obtaining second-order, condition-free algorithms for linear programming and $\ell_p$ regression \citep{bcll18,ls19,akps19,jls21}. Observe that $\ell_p$ regression is a special case of the problem we study in objective \eqref{eq:interpolation}, which is recovered when all $n_i=1$, and $\ell_{\infty}$ regression is captured by linear programming. Note that neither of these problem families is expressive enough to capture the objectives we study. In general, to achieve iteration complexities in the smaller of the two dimensions for these problems, it appears that a geometric understanding of the solution space is required --- these ideas were central to the improvements obtained by \citet{ls19,jls21} as well as our work.

\paragraph{Open problems.} 
Our work raises several open questions. One limitation of \Cref{mainthm:fair_regression_iteration_complexity} is that its iteration complexity is not high-accuracy, meaning its dependence on \( \eps \) is not \(\mathrm{polylog}(1/\eps)\). Designing a high-accuracy solver under the same conditions as \Cref{mainthm:gp_regression_iteration_complexity} with iteration complexity \(\widetilde{O}\inparen{\mathsf{poly}(\min\inbraces{\rank{\mA}, m}, \logv{\frac{1}{\eps}})}\) remains an open problem.

A more ambitious general goal is to design algorithms for convex quadratic programs with the aforementioned iteration complexity. This would generalize analogous results for linear programming \citep{ls19}. We view the current work as a first step towards this goal, as the objective \eqref{eq:main_objective} is a structured convex quadratic program for which we get an iteration complexity independent of $m$. It would also be interesting to consider other complexity measures beyond $\rank{\mA}$, for instance, assumptions about the ground-truth labeling vector $\vx_i^\star$ for each group's data $S_i$.

Finally, our results suggest that optimizing for ``$\ell_p$-interpolants'' between non-robust and robust objectives may be computationally easier than optimizing for the robust objective alone. A more precise statistical characterization of how robustness and utility trade-off as \( p \) varies in collaborative, fair, or multi-distributional learning settings would be valuable. Additionally, exploring interpolations or solution concepts along the Pareto frontier of the \( m \)-dimensional multi-objective optimization problem or other DRO notions (eg Wasserstein DRO \citep{bkm19, cpo20}) could yield further insights.


\subsection{Paper Outline}

In the remainder of this paper, we will outline the key details of our approach and provide a proof outline for our theoretical results. In \Cref{sec:gp_reg_overview}, we give proof sketches of our main results. In \Cref{sec:blw_to_l2}, we prove some background results that appear in the main body, particularly about block Lewis weights. In \Cref{sec:mirror_descent}, we give an analysis of mirror descent under inexact subproblem solves -- we will need this in the proof of \Cref{mainthm:gp_regression_iteration_complexity}. In \Cref{sec:optimal_ms_acceleration}, we modify an acceleration scheme due to \cite{chjjs22}, which we will use to iterate calls to the proximal subproblem solver \eqref{eq:gp_prox_intro_reg} for the proof of \Cref{mainthm:gp_regression_iteration_complexity}. In \Cref{sec:gp_robust_proof}, we prove \Cref{mainthm:fair_regression_iteration_complexity}. In \Cref{sec:gp_interpolating_proof}, we prove \Cref{mainthm:gp_regression_iteration_complexity}. Finally, in \Cref{app:exps} we include an empirical comparison of our proposed algorithms against the aforementioned baselines.